\documentclass[11pt]{article}
\usepackage[T1]{fontenc}
\usepackage{amsmath,amssymb,amsthm,mathtools,booktabs}
\usepackage[margin=1in]{geometry}
\usepackage[colorlinks=true,urlcolor=blue,citecolor=blue]{hyperref}
\newtheorem{theorem}{Theorem}
\newtheorem{lemma}[theorem]{Lemma}
\newtheorem{corollary}[theorem]{Corollary}
\newcommand{\rem}{\operatorname{rem}}
\newcommand{\calls}{\operatorname{calls}}
\newcommand{\eq}{\operatorname{eq}}
\title{Short remainder-only certificates for the greatest common divisor}
\author{}
\date{October 16, 2026}
\begin{document}
\maketitle
\begin{abstract}
We give a remainder-only straight-line derivation of the greatest common divisor
of two positive integers using $O(\log a/\log\log a)$ remainder calls, uniformly
for $a>b\geq1$. Previously generated values may be retained and reused, but
numerical constants and other arithmetic operations are not available as
instructions. The construction combines a stored Euclidean prefix with
three-call contractions and two smaller recursive derivations. It yields the
same upper bound for intrinsic remainder-call complexity, and therefore a
negative answer to the uniform-process strengthening discussed by Moschovakis.
It does not establish a deterministic remainder-only algorithm with this cost.
The mathematical proof below and the exact scope of the accompanying Lean
verification are distinguished in the final section.
\end{abstract}

\section{Statement and relation to the book}
Throughout, logarithms are to base $2$. A derivation starts with its two inputs.
An instruction adds $\rem(x,y)$, where $x,y$ are already available and $y>0$.
Earlier values remain available. No constants are supplied.
Let $G(a,b)$ be the minimum number of instructions needed to make $\gcd(a,b)$
available, and let $D(a,b)=G(a,b)$ for coprime inputs. Thus $G(a,0)=0$ and
$D(a,1)=0$. These are costs of generating values.

\begin{theorem}\label{thm:main}
Uniformly for integers $a>b\geq1$,
\[
G(a,b)=O\left(\frac{\log a}{\log\log a}\right).
\]
In particular this bound holds for $D(a,b)$ when $\gcd(a,b)=1$.
\end{theorem}
As usual, the asymptotic assertion concerns sufficiently large $a$; a bound
valid at every positive input is $O(1+\log(2+a)/\log\log(4+a))$.

Moschovakis's Main Conjecture concerns logarithmic lower bounds for gcd
algorithms from $\rem,\eq_0$ \cite[p.~2]{moschovakis2019}. The uniform-process
strengthening is discussed after Corollary~6C.7 \cite[p.~188]{moschovakis2019};
see also (4F-5) on p.~151. Our conclusion addresses that strengthening.
The definitions of generated certificates and intrinsic norms are in
\cite[\S\S1D, 4C--4E]{moschovakis2019}. Theorem~4F.1
\cite[pp.~149--150]{moschovakis2019} transfers pointwise intrinsic bounds to a
single uniform process. Deterministic uniform processes are treated separately
in \cite[\S4H]{moschovakis2019}.

\section{Normalization}
Write $g=\gcd(a,b)>0$. For positive $y$,
\[
\rem(gx,gy)=g\rem(x,y).
\]
Every generated value is divisible by $g$. Consequently scaling a derivation,
or dividing all its values by $g$, preserves its length and the validity of
every instruction. Hence
\begin{equation}\label{eq:normalize}
G(a,b)=D(a/g,b/g).
\end{equation}
The normalized pair $(1,0)$ is understood to have cost zero. Define
\[
T(n)=\max_{0\leq b\leq a\leq 2^n,\ a>0}G(a,b).
\]
Euclid's algorithm makes this a finite maximum. Normalization reduces an upper
bound for $T$ to coprime inputs, with equality, zero inputs, and inputs equal to
one handled without calls.

\section{Three-call contractions}
\begin{lemma}\label{lem:gap}
If $m>0$ is available and a finite collection $S$ of available values occupies
$t\geq2$ residue classes modulo $m$, at most three remainder calls generate $w$
with
\[
0<w<\frac{m}{t-1}.
\]
\end{lemma}
\begin{proof}
Sort the distinct residues. Some adjacent pair $u<v$ has gap
$\delta=v-u<m/(t-1)$, because the sum of the $t-1$ adjacent gaps is less than
$m$. Each selected residue can be generated with one call from a witness in
$S$. If $u=0$, use $v$. If $0<u\leq\delta$, use $u$. Otherwise
$0<\delta<u$ and $v=u+\delta$, so $\rem(v,u)=\delta$. This uses at most three
calls. Subtraction and sorting describe the proof and the selection of
instructions; they are not instructions of the derivation.
\end{proof}

\begin{lemma}[Euclidean-prefix lemma]\label{lem:prefix}
Let $r_0=a>r_1=b>\cdots>r_h>0$ be a Euclidean prefix of a coprime pair, with
$r_{i-1}=q_i r_i+r_{i+1}$ and $q_i\geq1$. For every integer $m>2a/r_h$, the
$h$ pairs
\[
(r_i\bmod m,r_{i+1}\bmod m),\qquad 0\leq i<h,
\]
are distinct. The prefix therefore occupies $t$ residue classes with $t^2\geq h$.
\end{lemma}
\begin{proof}
Suppose the pairs at $i<j$ coincide. Set
\[
M=\prod_{\ell=i+1}^{j}\begin{pmatrix}q_\ell&1\\1&0\end{pmatrix}
 =\begin{pmatrix}A&B\\C&E\end{pmatrix},\qquad
\binom{r_i}{r_{i+1}}=M\binom{r_j}{r_{j+1}}.
\]
All entries are nonnegative, $E\leq A$, and $\det M=(-1)^{j-i}$.
Also $Ar_j\leq r_i\leq a$, whence $A\leq a/r_h$.
Write
\[
\Delta=\det(M-I)=1+\det M-A-E.
\]
For an odd number of factors, $\Delta=-A-E<0$. A product of at least two
factors has $A\geq2$ and $E\geq1$; thus for an even positive number of factors,
$\Delta=2-A-E<0$. In either case,
\[
0<|\Delta|\leq A+E\leq2A\leq2a/r_h<m.
\]
The presumed congruence gives $(M-I)(r_j,r_{j+1})^{\mathsf T}\equiv0\pmod m$.
Multiplication by the adjugate shows that $m$ divides both $\Delta r_j$ and
$\Delta r_{j+1}$. Consecutive remainders are coprime, so B\'ezout's identity
implies $m\mid\Delta$, a contradiction. Finally, $t$ individual residue classes
give at most $t^2$ ordered pairs.
\end{proof}
Thus a stored prefix can be reused for every available modulus $m>2a/r_h$.
When $\sqrt h>1$, three calls suffice to generate
\begin{equation}\label{eq:contract}
0<w<\frac{m}{\sqrt h-1}.
\end{equation}
The universal quantifier over $m$ is essential: the modulus generated by a
contraction is not predetermined.

\section{Construction and recurrence}
Fix a sufficiently large integer $n$ and put $N=2^n$ and $h=\lceil\sqrt n\rceil$.
After normalization we may assume $a>b>1$ and $a\leq N$.
Generate the Euclidean prefix until its last positive value is at most
$N^{3/4}$ or until $r_h$ is available. Check the input $b$ as part of this
rule. The preliminary work costs at most $h$ calls. Euclid cannot encounter
zero first without having encountered the positive remainder $1$.

If a positive remainder $c\leq N^{3/4}$ is encountered, generate the next
remainder $s<c$ with one more call. The pair $(c,s)$ is coprime. If $s=0$, then
$c=1$ is already available; otherwise both inputs to the remaining derivation
are at most $2^{\lceil3n/4\rceil}$. This gives the bound
\begin{equation}\label{eq:early}
h+1+T(\lceil3n/4\rceil).
\end{equation}

Otherwise $r_h>N^{3/4}$, so $2a/r_h<2N^{1/4}$. For $n>12$,
$N^{1/3}>2N^{1/4}$, and Lemma~\ref{lem:prefix} applies to every available
$m>N^{1/3}$. Starting at $m=r_h$, repeatedly use the contraction inequality until
$0<m\leq N^{1/3}$. Set
\[
\kappa_n=\sqrt h-1>1,\qquad
J_n=\left\lceil\frac{2n/3}{\log\kappa_n}\right\rceil.
\]
There are at most $J_n$ contractions: if $J_n$ have occurred then
$m<N/\kappa_n^{J_n}\leq N^{1/3}$. Since
$\log\kappa_n=\tfrac14\log n+O(1)$, their cost is $3J_n=O(n/\log n)$.

The small $m$ need not be coprime to either input. Generate
$u=\rem(a,m)$ and $v=\rem(b,m)$, then concatenate derivations of
$d=\gcd(m,u)$ and $e=\gcd(d,v)$. These derivations exist by the definition of
$T$ and use only inputs bounded by $m$. They are not gcd-oracle instructions.
The gcd identities give
\[
e=\gcd(\gcd(m,a\bmod m),b\bmod m)
 =\gcd(m,\gcd(a,b))=1.
\]
The cost in this case is at most
\begin{equation}\label{eq:late}
h+3J_n+2+2T(\lceil n/3\rceil).
\end{equation}
Retaining values makes concatenation legitimate. In particular, the second
recursive derivation has both $d$ and $v$ available even if the first one uses
neither original input.

The two case estimates and normalization imply
\begin{equation}\label{eq:recurrence}
T(n)\leq\max\{T(\lceil3n/4\rceil)+O(\sqrt n),\;
2T(\lceil n/3\rceil)+O(n/\log n)\}.
\end{equation}
All implied constants are independent of the inputs.

\begin{proof}[Proof of Theorem~\ref{thm:main}]
For $n\geq2$ let $f(n)=n/\log n$. The ratios
\[
\frac{f(\lceil3n/4\rceil)}{f(n)}\longrightarrow\frac34,
\qquad
\frac{2f(\lceil n/3\rceil)}{f(n)}\longrightarrow\frac23
\]
show that, beyond a fixed threshold, the two recursive terms have potentials
at most $\tfrac78 f(n)$ and $\tfrac56 f(n)$, respectively.
Both overheads are at most $Af(n)$ for a fixed $A$. Choose $C\geq8A$ and
increase $C$ to cover the finitely many initial values with indices at least
$2$. Choose the induction threshold large enough that both recursive indices
are at least $2$ and strictly less than $n$.
Strong induction applied to this recurrence then gives $T(n)\leq Cf(n)$:
the two bounds are $(7C/8+A)f(n)$ and $(5C/6+A)f(n)$, both at most $Cf(n)$.
Finally take $n=\lceil\log a\rceil$. This also gives $G(a,b)=o(\log a)$
uniformly as $a\to\infty$.
\end{proof}

\section{Certificates and uniform processes}
Use the book's partial remainder operation: its second argument must be
nonzero. Write $\mathcal N_0=(\mathbb N,\rem,\eq_0)$ and
$\mathcal N_\varepsilon=(\mathbb N,\rem,\eq_0,\eq_1)$; compare
\cite[(1C-10), pp.~24, 142--143]{moschovakis2019}.
A derivation determines a finite generated partial substructure by recording
its remainder equations. Repeated equations can be discarded.

For the relation $\gcd(a,b)=1$ on all natural inputs, a derivation of $1$
together with $\eq_1(1)=\mathrm{tt}$ certifies the relation: every common
divisor of the images of the inputs divides every transported generated
value, hence divides $1$. Conversely, without a positive $\eq_1$ fact the
map $x\mapsto2x$ preserves every recorded remainder and zero-test fact, as
well as every negative one-test fact. It cannot preserve coprimeness. A
certificate therefore reaches $1$. Enumerating a finite generated diagram
in generation order gives a derivation using no more calls than the number
of remainder facts. Consequently, on coprime inputs,
\[
\calls_{\rem}(\mathcal N_\varepsilon,\gcd=1,a,b)=D(a,b),\qquad
\calls(\mathcal N_\varepsilon,\gcd=1,a,b)=D(a,b)+1.
\]
If the relation is instead defined as $a>0\wedge b>0\wedge\gcd(a,b)=1$,
add the two input-positivity facts. The remainder count is unchanged, but
one should state the full count as $D(a,b)+O(1)$ rather than assert the exact
$D+1$ formula without addressing positivity. For example, the single fact
$\eq_1(1)=\mathrm{tt}$ at $(2,1)$ allows the image $(0,1)$.

For gcd as a function, the certificate must reach its output $g$, giving a
lower bound of $G(a,b)$. For an upper bound, append
\[
\rem(a,g)=z,\qquad \rem(b,g)=z,\qquad \eq_0(z)=\mathrm{tt},
\]
where the generated $z$ is zero. Every homomorphism $\pi$ has $\pi(z)=0$ and
$\pi(g)>0$ (by definedness of either appended remainder). Thus $\pi(g)$
divides both input images. Conversely every common divisor of those images
divides $\pi(g)$ by induction over the derivation. Therefore
\[
G(a,b)\leq\calls_{\rem}(\mathcal N_0,\gcd,a,b)\leq G(a,b)+2.
\]
Positivity facts may be appended if the function is restricted to positive
inputs. For a noncoprime input, add $\eq_1(g)=\mathrm{ff}$ to this gcd
certificate to certify noncoprimeness. Thus the sublogarithmic remainder-call
bound holds for both truth values of coprimeness as well as for gcd.
Theorem~4F.1 supplies a single uniform process realizing each intrinsic norm
pointwise \cite[Theorem~4F.1]{moschovakis2019}.

\section{Scope and verification}
Selecting a close pair of residues is uncharged existential choice in a
straight-line derivation. A deterministic remainder-only algorithm cannot
assume that choice is available at constant cost. Computing all prefix
residues would take about $h$ calls per contraction. The construction therefore
does not settle the deterministic conjecture.

The supplied offline Python planner uses unrestricted arithmetic to select
instructions; its independent replay checks operand availability, nonzero
divisors, every remainder equation, and the final gcd. Running its default
seed reproduces:
\begin{center}
\begin{tabular}{rrr}\toprule
Input bits & Derivation calls & Euclid calls to reach $1$\\\midrule
512 & 196 & 326\\
2048 & 669 & 1204\\
8192 & 2059 & 4754\\\bottomrule
\end{tabular}
\end{center}
These are checks of finite instances, not an asymptotic proof. A newly
specified reproducible suite also passes 310 cases, including 175 noncoprime
pairs, with 3,543 contractions across 335 reservoir phases. Four deliberately
invalid traces are rejected. This suite is recorded in the repository; it is
not claimed to reconstruct the unspecified original 310-case suite.

\subsection{Machine-checked statements}
The Lean 4.35.0 development, with mathlib pinned to version 4.35.0, proves the
uniform real-logarithmic bound, its worst-case Big-O formulation, the
Euclidean-prefix lemma, the complete three-call contraction lemma, normalization,
and the certificate upper bounds. It contains no proof placeholders or added
axioms. An axiom audit of the main results reports only the standard Lean
foundations \texttt{propext}, \texttt{Classical.choice}, and \texttt{Quot.sound}.
The theorem \texttt{ShortGCD.uniform\_real\_bound} states
\[
\exists C>0\ \forall a,b\in\mathbb N,\quad
4\leq a\ \wedge\ b\leq a\ \Longrightarrow\
G(a,b)\leq C\frac{\log_2 a}{\log_2(\log_2 a)}.
\]
The proof starts from a finite-store instruction semantics in which the only
instruction is remainder with a positive, available divisor. Costs are indices
of derivations, rather than an assumed recurrence. A formal conversion to
finite generated diagrams proves the gcd certificate bound with two extra
remainder facts and the bound for positive and negative coprimeness instances.
The equality with the existing $D$ definition is also checked.

The intrinsic complexities themselves are also defined in Lean, directly from
certificates: \texttt{callsCoprimeness} for the coprimeness relation in
$N_\varepsilon$ and \texttt{callsGcd} for the gcd without $\eq_1$ facts. On
coprime pairs the former agrees with the existing reduction to $D$
(\texttt{callsCoprimeness\_eq\_callsRem}). The theorems
\texttt{ShortGCD.not\_conjecture\_3\_4} and
\texttt{ShortGCD.not\_main\_conjecture\_4F5} state that no $r>0$ satisfies
$\mathrm{calls}(\mathrm{rem})(N_\varepsilon,\perp\!\!\!\perp,a,b)\geq r\log_2 a$
for infinitely many $a\geq b\geq1$, and that no $K>0$ satisfies
$2\log_2\max(x,y)\leq K\,\mathrm{calls}(\mathrm{rem})(\gcd,x,y)$ for infinitely
many pairs. The general theory of uniform processes and Theorem~4F.1 are cited
from the book, not re-formalized here. They are needed only to interpret these
intrinsic bounds as statements about a single best uniform process, not for the
refutations themselves.

\subsection{Integer parameters used in Lean}
To avoid real rounding in the construction itself, the formal proof uses, for
$j=\lfloor\log_2 n\rfloor\geq32$,
\[
p=\lfloor j/4\rfloor,\quad k=2^p,\quad h=2^{2p}+1,\quad
J=\lfloor n/p\rfloor+1,\quad
\ell=\lceil3n/4\rceil,\quad s=\lceil n/3\rceil.
\]
The finite construction proves, for arbitrary natural parameters satisfying
$k\geq1$, $k^2<h$, $L\geq1$, $2A\leq ML$, and $A\leq Mk^J$,
\[
S(A)\leq\max\{h+1+S(L),\ h+3J+2+2S(M)\},
\qquad S(A)=\max_{a,b\leq A}G(a,b).
\]
Taking $A=2^n,L=2^\ell,M=2^s$ gives the required hypotheses. Lean checks
\[
(h+1)j\leq64n,\quad (h+3J+2)j\leq64n,\quad
j-2\leq\log_2^{\mathrm{nat}}\ell,\log_2^{\mathrm{nat}}s,
\quad 8\ell\leq7n,\quad16s\leq7n.
\]
The last estimates and $8j\leq9(j-2)$ give a factor $63/64$ for either
recursive contribution after clearing denominators. Taking $C\geq4096$ and
large enough for the finite initial segment proves
$S(2^n)\lfloor\log_2 n\rfloor\leq Cn$ for all $n$.
Comparing integer floors and ceilings with real base-two logarithms yields the
stated real bound. These parameter choices change only absolute constants,
not the proof mechanism or its asymptotic order.

\begin{thebibliography}{1}
\bibitem{moschovakis2019}
Yiannis N. Moschovakis.
\emph{Abstract Recursion and Intrinsic Complexity}.
Lecture Notes in Logic, vol.~48. Cambridge University Press and Association
for Symbolic Logic, 2019.
\href{https://www.math.ucla.edu/~ynm/lectures/aricplace/c-aric.pdf}{Author's
prepublication PDF}, compiled November 30, 2018.
Page references above are to printed page numbers, not PDF page indices.
\end{thebibliography}
\end{document}
